\documentclass[11pt,letterpaper]{article}

% PAGE FORMAT
\usepackage[
    letterpaper,
    top=3cm,
    botom=3cm,
    left=2cm,
    right=2cm
]{geometry}

% ENCODING AND TYPOGRAPHY
\usepackage[T1]{fontenc}
\usepackage[utf8]{inputenc}
\usepackage{parskip}

% MATHEMATICS
\usepackage{amsmath}
\usepackage{amsfonts}
\usepackage{amssymb}
\usepackage{amsthm}

% Tables
\usepackage{array}
\usepackage{longtable}
\usepackage{booktabs}
\usepackage{colortbl}
\usepackage[table]{xcolor}
\usepackage{multirow}

\newtheorem{theorem}{Theorem}[section]
\newtheorem{proposition}[theorem]{Proposition}
\newtheorem{definition}[theorem]{Definition}
\newtheorem{remark}[theorem]{Remark}

% FIGURES AND ALGORITHMS
\usepackage{graphicx}
\usepackage{caption}
\usepackage{float}
\usepackage{algorithm}
\usepackage{algorithmic}

% BIBLIOGRAPHY
\usepackage[
    backend=biber,
    style=ieee
]{biblatex}
\addbibresource{references.bib}
\addbibresource{references_modelo.bib}

\usepackage{hyperref}

% ---- colores del cronograma ----
\definecolor{c1fill}{RGB}{232,230,252}
\definecolor{c1acc}{RGB}{83,74,183}
\definecolor{c2fill}{RGB}{220,245,235}
\definecolor{c2acc}{RGB}{15,110,86}
\definecolor{c3fill}{RGB}{220,234,248}
\definecolor{c3acc}{RGB}{24,95,165}
\definecolor{c4fill}{RGB}{254,242,220}
\definecolor{c4acc}{RGB}{133,79,11}
\definecolor{entfill}{RGB}{253,235,235}
\definecolor{entacc}{RGB}{163,45,45}
\definecolor{headbg}{RGB}{45,45,50}

% TITLE
\title{
    \textbf{\LARGE Anteproyecto: Influencia de la Configuración Geométrica
    de una Telaraña Orbicular en su Dinámica Vibratoria
    Mediante un Modelo de Mecánica Analítica}
}
\author{
    Amelia Hoyos Vélez,
    Nicolle Rodríguez Estupiñán,
    Miguel Ángel León Botero\\[0.3cm]
    \small Ingeniería Física,
    Escuela de Ciencias e Ingeniería,
    Universidad EAFIT
}
\date{2026}

\begin{document}
\maketitle


\section{Justificación}

Las arañas orbitelas construyen estructuras cuya geometría no es azarosa:
el número de hilos radiales, el espaciado entre los anillos espirales y la
tensión impuesta en cada elemento son parámetros que la araña puede modificar
activamente \cite{mortimer2016,mortimer2019}. Esta capacidad sugiere que la
arquitectura de la red cumple una función sensorial específica, pues la araña
detecta posición, dirección y magnitud de perturbaciones externas a través de
los patrones vibratorios que se propagan por los hilos hasta sus patas
\cite{masters1981,mortimer2019}. Sin embargo, no toda configuración geométrica
favorece la transmisión de información relevante: una red inadecuada puede
atenuar señales de presa y amplificar ruido ambiental ---viento, contacto con
hojas--- comprometiendo la capacidad de caza \cite{mortimer2016,nakata2010}.

Esto plantea una pregunta fundamental de mecánica estructural aplicada a
sistemas biológicos:

\medskip
\begin{center}
\textit{¿Cómo determina la configuración geométrica de una telaraña orbicular
las propiedades vibratorias de la estructura y la forma en que una perturbación
mecánica se transmite entre sus nodos?}
\end{center}
\medskip

Responder esta pregunta requiere un modelo dinámico que relacione
explícitamente los parámetros geométricos de la red con su espectro de
frecuencias naturales, sus modos de vibración y su respuesta temporal ante
perturbaciones. El formalismo lagrangiano ofrece el marco adecuado para
construir dicho modelo de manera sistemática.

\section{Objetivo General}
Analizar la influencia de la configuración geométrica de una telaraña orbicular en su dinámica vibratoria, mediante el modelado del sistema y el estudio de su respuesta ante perturbaciones mecánicas, con el fin de determinar cómo la distribución de sus nodos e hilos, así como la tensión generada entre ellos, afectan las propiedades dinámicas de la estructura.

\section{Objetivos Específicos}
\begin{itemize}
    \item Modelar la dinámica de una telaraña orbicular considerando su distribución geométrica, mediante la formulación de las ecuaciones de movimiento del sistema a partir de la mecánica analítica para establecer las bases matemáticas de su análisis vibratorio.

    \item Caracterizar el comportamiento vibratorio del sistema ante una perturbación mecánica mediante el análisis de sus frecuencias naturales, modos normales y respuesta dinámica, para determinar cómo se transmite la perturbación a través de la estructura.

    \item Simular numéricamente la respuesta vibratoria de la telaraña ante diferentes configuraciones geométricas y perturbaciones mecánicas, para contrastar los resultados obtenidos con el análisis teórico.
\end{itemize}

\section{Planteamiento del Problema}

 

\subsection{Geometría de la red y posición de equilibrio}

Se considera una telaraña orbicular tridimensional compuesta por $N$ hilos
radiales y $M$ anillos espirales \cite{aoyanagi2010,lott2022}. En equilibrio,
la telaraña yace en el plano $z=0$. Los nodos de la red se etiquetan con el
par de índices $(n,\,m)$, donde

\begin{equation}
    n = 0, 1, \dots, N-1 \quad \text{(índice radial)}, \qquad
    m = 0, 1, \dots, M \quad \text{(índice de anillo)}.
\end{equation}

El caso $m = 0$ corresponde al nodo central, que se supone fijo
($\mathbf{X}^0_{n,0} = \mathbf{0}$,
$\mathbf{u}_{n,0} \equiv \mathbf{0}$ para todo $n$). Los nodos del borde
exterior ($m = M$) son libres. Para $m \geq 1$, la posición de equilibrio
del nodo $(n,m)$ es

\begin{equation}\label{eq:equil}
    \mathbf{X}^{0}_{n,m}
    = \begin{pmatrix} r_m\cos\theta_n \\ r_m\sin\theta_n \\ 0
      \end{pmatrix},
    \qquad \theta_n = \frac{2\pi n}{N},
\end{equation}

donde $r_m > 0$ es el radio del $m$-ésimo anillo, con
$r_1 < r_2 < \cdots < r_M$.

\subsection{Desplazamientos y coordenadas generalizadas}

Ante una perturbación, cada nodo libre $(n,m)$ se desplaza en las tres
direcciones espaciales. Se define el vector de desplazamiento

\begin{equation}\label{eq:despl}
    \mathbf{u}_{n,m}(t)
    = \mathbf{x}_{n,m}(t) - \mathbf{X}^{0}_{n,m}
    = \begin{pmatrix} u^x_{n,m}(t) \\ u^y_{n,m}(t) \\ u^z_{n,m}(t)
      \end{pmatrix} \in \mathbb{R}^3,
\end{equation}

donde $(u^x_{n,m},\, u^y_{n,m})$ son los desplazamientos en el plano
(modos longitudinales) y $u^z_{n,m}$ es el desplazamiento transversal
(modos fuera del plano). Reuniendo todos los desplazamientos en un único
vector de estado, ordenado primero por anillo y luego por radio,

\begin{equation}\label{eq:q}
    \mathbf{q}(t)
    = \bigl(
        \mathbf{u}_{0,1}^\top,\,
        \mathbf{u}_{1,1}^\top,\,\dots,\,
        \mathbf{u}_{N-1,1}^\top,\;\dots,\;
        \mathbf{u}_{0,M}^\top,\,\dots,\,
        \mathbf{u}_{N-1,M}^\top
      \bigr)^\top
    \in \mathbb{R}^{3NM}.
\end{equation}

El sistema tiene $3NM$ grados de libertad activos.

\subsection{Conectividad de la red}

La red presenta dos tipos de hilo, cada uno caracterizado por su constante
elástica y su tensión inicial (pretensión) \cite{aoyanagi2010,soler2016}:

\begin{itemize}
    \item \emph{Hilo radial} $(n,m)$: conecta el nodo $(n,m-1)$ con el
    nodo $(n,m)$, con constante elástica $K_m > 0$ y pretensión
    $T^R_m > 0$.
    \item \emph{Hilo espiral} $(n,m)$: conecta el nodo $(n-1,\,m)$ con el
    nodo $(n,m)$, con constante elástica $k_m > 0$ y pretensión
    $T^S_m > 0$ (índice $n$ en módulo $N$).
\end{itemize}

Para un hilo genérico que une los nodos $i$ y $j$, se definen la longitud
de equilibrio y el vector unitario de equilibrio:

\begin{equation}\label{eq:L0}
    L^0_{ij} = \lVert \mathbf{X}^0_i - \mathbf{X}^0_j \rVert,
    \qquad
    \hat{\mathbf{e}}^0_{ij}
    = \frac{\mathbf{X}^0_i - \mathbf{X}^0_j}{L^0_{ij}}
    \in \mathbb{R}^3.
\end{equation}

\subsection{Energía cinética}

Modelando la masa de los hilos como concentrada en los nodos \cite{meirovitch2001},
la energía cinética del sistema es

\begin{equation}\label{eq:T}
    T = \frac{1}{2}
        \sum_{n=0}^{N-1}\sum_{m=1}^{M}
        \mu_{n,m}\,\dot{\mathbf{u}}_{n,m}^{\top}\dot{\mathbf{u}}_{n,m},
\end{equation}

donde $\mu_{n,m} > 0$ es la masa concentrada en el nodo $(n,m)$.

\subsection{Elongación de un hilo y energía potencial}

Cuando los nodos $i$ y $j$ se desplazan, la elongación del hilo que los
une respecto a su longitud de equilibrio se aproxima, para desplazamientos
pequeños ($\lVert\mathbf{u}_i - \mathbf{u}_j\rVert \ll L^0_{ij}$),
mediante la proyección del desplazamiento relativo sobre el vector unitario
de equilibrio \cite{goldstein2002,meirovitch2001}:

\begin{equation}\label{eq:elong}
    \Delta l_{ij}
    = L_{ij} - L^0_{ij}
    \approx \hat{\mathbf{e}}^0_{ij}\cdot(\mathbf{u}_i - \mathbf{u}_j).
\end{equation}

La energía potencial de cada hilo tiene dos contribuciones. La
\emph{contribución elástica} proviene de la deformación axial:

\begin{equation}\label{eq:Velast}
    V^{\text{elás}}_{ij}
    = \frac{1}{2}\kappa_{ij}(\Delta l_{ij})^2
    = \frac{1}{2}\kappa_{ij}
      \bigl[\hat{\mathbf{e}}^0_{ij}\cdot(\mathbf{u}_i-\mathbf{u}_j)\bigr]^2.
\end{equation}

La \emph{contribución geométrica} proviene de la pretensión $T^0_{ij}$:
un hilo pretensado resiste el movimiento transversal a su eje incluso sin
elongación neta \cite{jayaraman1981,wirth2011}. Expandiendo la longitud instantánea a segundo orden en
el desplazamiento relativo $\boldsymbol{\delta}_{ij}=\mathbf{u}_i-\mathbf{u}_j$,

\begin{equation}
    L_{ij}
    \approx L^0_{ij}
    + \hat{\mathbf{e}}^0_{ij}\cdot\boldsymbol{\delta}_{ij}
    + \frac{1}{2L^0_{ij}}
      \Bigl[
          \lVert\boldsymbol{\delta}_{ij}\rVert^2
          - \bigl(\hat{\mathbf{e}}^0_{ij}\cdot\boldsymbol{\delta}_{ij}\bigr)^2
      \Bigr],
\end{equation}

el término entre corchetes corresponde a la componente de
$\boldsymbol{\delta}_{ij}$ perpendicular a $\hat{\mathbf{e}}^0_{ij}$,
cuyo cuadrado genera la energía geométrica:

\begin{equation}\label{eq:Vgeom}
    V^{\text{geom}}_{ij}
    = \frac{T^0_{ij}}{2L^0_{ij}}
      \Bigl[
          \lVert\boldsymbol{\delta}_{ij}\rVert^2
          - \bigl(\hat{\mathbf{e}}^0_{ij}\cdot\boldsymbol{\delta}_{ij}\bigr)^2
      \Bigr].
\end{equation}

Sumando ambas contribuciones sobre todos los hilos, la energía potencial
total es

\begin{equation}\label{eq:V}
    V = \frac{1}{2}\sum_{(i,j)}
    \left\{
        \kappa_{ij}
        \bigl[\hat{\mathbf{e}}^0_{ij}\cdot(\mathbf{u}_i-\mathbf{u}_j)\bigr]^2
        +
        \frac{T^0_{ij}}{L^0_{ij}}
        \Bigl[
            \lVert\mathbf{u}_i-\mathbf{u}_j\rVert^2
            -\bigl(\hat{\mathbf{e}}^0_{ij}
            \cdot(\mathbf{u}_i-\mathbf{u}_j)\bigr)^2
        \Bigr]
    \right\}.
\end{equation}

\subsection{Ecuaciones de movimiento}

El lagrangiano del sistema es $\mathcal{L} = T - V$. Aplicando las
ecuaciones de Euler--Lagrange a cada coordenada generalizada
$\mathbf{u}_{n,m}$,

\begin{equation}\label{eq:EL}
    \frac{\mathrm{d}}{\mathrm{d}t}
    \frac{\partial\mathcal{L}}{\partial\dot{\mathbf{u}}_{n,m}}
    -
    \frac{\partial\mathcal{L}}{\partial\mathbf{u}_{n,m}}
    = \mathbf{0},
\end{equation}

se obtiene la ecuación de movimiento para cada nodo libre $(n,m)$:

\begin{equation}\label{eq:eom_local}
    \mu_{n,m}\,\ddot{\mathbf{u}}_{n,m}
    =
    -\sum_{j\in\mathcal{N}(n,m)}
    \mathbf{k}^e_{(n,m),j}
    \bigl(\mathbf{u}_{n,m} - \mathbf{u}_j\bigr),
\end{equation}

donde $\mathcal{N}(n,m)$ es el conjunto de nodos adyacentes a $(n,m)$ y
$\mathbf{k}^e_{ij} \in \mathbb{R}^{3\times 3}$ es la \emph{matriz de
rigidez elemental} del hilo $(i,j)$, obtenida al derivar la energía
potencial \eqref{eq:V} \cite{jayaraman1981,meirovitch2001}:

\begin{equation}\label{eq:ke}
    \mathbf{k}^e_{ij}
    =
    \frac{T^0_{ij}}{L^0_{ij}}\,\mathbf{I}_3
    +
    \left(\kappa_{ij} - \frac{T^0_{ij}}{L^0_{ij}}\right)
    \hat{\mathbf{e}}^0_{ij}\bigl(\hat{\mathbf{e}}^0_{ij}\bigr)^\top.
\end{equation}

El primer término de \eqref{eq:ke} es la \emph{rigidez geométrica},
proporcional a la pretensión, que actúa isotrópicamente sobre las tres
componentes del desplazamiento relativo, incluida la dirección transversal
al plano. El segundo término es la \emph{rigidez elástica adicional},
proyectada sobre la dirección axial del hilo. Sin pretensión
($T^0_{ij}=0$), los modos transversales tendrían frecuencia nula; la
pretensión los dota de frecuencias finitas determinadas por la relación
$T^0_{ij}/L^0_{ij}$ \cite{lott2022,frohlich1982,wirth2011}.

\subsection{Construcción de las matrices globales}

Reuniendo las ecuaciones \eqref{eq:eom_local} para todos los nodos libres
en el vector de estado $\mathbf{q}$ definido en \eqref{eq:q}, el sistema
adopta la forma matricial

\begin{equation}\label{eq:sistema}
    \mathbf{M}\,\ddot{\mathbf{q}} + \mathbf{K}\,\mathbf{q} = \mathbf{f}(t),
\end{equation}

donde $\mathbf{f}(t) \in \mathbb{R}^{3NM}$ es el vector de forzamiento
externo.

\paragraph{Matriz de masa.}
La energía cinética \eqref{eq:T} es diagonal en $\mathbf{q}$, por lo que
$\mathbf{M}$ es una matriz diagonal de bloques $3\times 3$:

\begin{equation}\label{eq:M}
    \mathbf{M}
    =
    \mathrm{diag}\!\bigl(
        \mu_{0,1}\mathbf{I}_3,\;
        \mu_{1,1}\mathbf{I}_3,\;
        \dots,\;
        \mu_{N-1,M}\mathbf{I}_3
    \bigr)
    \in \mathbb{R}^{3NM\times 3NM}.
\end{equation}

\paragraph{Matriz de rigidez.}
Cada hilo $(i,j)$ aporta a la rigidez global mediante el bloque de
ensamblaje estándar \cite{meirovitch2001,jayaraman1981}

\begin{equation}\label{eq:ensamble}
    \Delta\mathbf{K}_{ij}
    =
    \begin{pmatrix}
         \mathbf{k}^e_{ij} & -\mathbf{k}^e_{ij} \\
        -\mathbf{k}^e_{ij} &  \mathbf{k}^e_{ij}
    \end{pmatrix}
    \in \mathbb{R}^{6\times 6},
\end{equation}

distribuido en las filas y columnas de $\mathbf{K}$ correspondientes a
los grados de libertad de los nodos $i$ y $j$. Cuando uno de los nodos
es el centro fijo, solo se ensambla el bloque $\mathbf{k}^e_{ij}$
correspondiente al nodo libre. La matriz global resulta de sumar sobre
todos los hilos:

\begin{equation}\label{eq:K}
    \mathbf{K}
    =
    \underbrace{
        \sum_{n=0}^{N-1}\sum_{m=1}^{M}
        \Delta\mathbf{K}^{R}_{(n,m-1),(n,m)}
    }_{\text{hilos radiales}}
    +
    \underbrace{
        \sum_{n=0}^{N-1}\sum_{m=1}^{M}
        \Delta\mathbf{K}^{S}_{(n-1,m),(n,m)}
    }_{\text{hilos espirales}},
\end{equation}

donde $\Delta\mathbf{K}^{R}$ usa $\kappa=K_m$, $T^0=T^R_m$ y
$\hat{\mathbf{e}}^0=\hat{\mathbf{e}}^{R}_{n,m}$, mientras que
$\Delta\mathbf{K}^{S}$ usa $\kappa=k_m$, $T^0=T^S_m$ y
$\hat{\mathbf{e}}^0=\hat{\mathbf{e}}^{S}_{n,m}$. La matriz $\mathbf{K}$
es simétrica, semidefinida positiva y rala de dimensión
$3NM\times 3NM$.

\subsection{Análisis modal}

Las frecuencias naturales $\omega_k$ y los modos normales
$\boldsymbol{\phi}_k$ se obtienen del problema de valores propios
generalizado:

\begin{equation}\label{eq:evp}
    \mathbf{K}\,\boldsymbol{\phi}_k
    = \omega_k^2\,\mathbf{M}\,\boldsymbol{\phi}_k,
    \qquad k = 1, 2, \dots, 3NM.
\end{equation}

Los modos satisfacen la ortogonalidad $\mathbf{M}$-normalizada:

\begin{equation}\label{eq:ortog}
    \boldsymbol{\phi}_k^\top\mathbf{M}\,\boldsymbol{\phi}_l = \delta_{kl},
    \qquad
    \boldsymbol{\phi}_k^\top\mathbf{K}\,\boldsymbol{\phi}_l
    = \omega_k^2\,\delta_{kl}.
\end{equation}

Los $3NM$ modos se separan naturalmente en dos familias: los \emph{modos
en el plano} (componentes $u^x$, $u^y$ dominantes), controlados
principalmente por las constantes elásticas $\kappa_{ij}$, y los
\emph{modos transversales} (componente $u^z$ dominante), cuya frecuencia
está determinada a orden dominante por el término de rigidez geométrica
$T^0_{ij}/L^0_{ij}$ \cite{lott2022,mortimer2019decoding,frohlich1982}.

\subsection{Respuesta temporal ante una perturbación}

Cuando una perturbación puntual se aplica sobre el nodo $(n^*,m^*)$ en
la dirección $\hat{\mathbf{d}}$, el vector de forzamiento es

\begin{equation}\label{eq:forzamiento}
    \mathbf{f}(t) = f(t)\,\mathbf{e}_{n^*,m^*,\hat{d}},
\end{equation}

donde $\mathbf{e}_{n^*,m^*,\hat{d}}$ es el vector canónico con
$\hat{\mathbf{d}}$ en las posiciones del nodo $(n^*,m^*)$ y cero en el
resto. Expandiendo la solución en la base modal,

\begin{equation}\label{eq:modal_exp}
    \mathbf{q}(t) = \sum_{k=1}^{3NM}\eta_k(t)\,\boldsymbol{\phi}_k,
\end{equation}

y usando la ortogonalidad \eqref{eq:ortog}, el sistema \eqref{eq:sistema}
se desacopla en $3NM$ osciladores independientes \cite{inman2014,meirovitch2001}:

\begin{equation}\label{eq:eta_k}
    \ddot{\eta}_k + \omega_k^2\,\eta_k = f_k(t),
    \qquad
    f_k(t) = \boldsymbol{\phi}_k^\top\mathbf{f}(t)
            = f(t)\,\boldsymbol{\phi}_k^\top\mathbf{e}_{n^*,m^*,\hat{d}}.
\end{equation}

La solución de cada coordenada modal mediante la fórmula de Duhamel es

\begin{equation}\label{eq:duhamel}
    \eta_k(t)
    = \eta_k(0)\cos(\omega_k t)
    + \frac{\dot{\eta}_k(0)}{\omega_k}\sin(\omega_k t)
    + \frac{1}{\omega_k}
      \int_0^t f_k(\tau)\sin\!\bigl(\omega_k(t-\tau)\bigr)\,\mathrm{d}\tau,
\end{equation}

con condiciones iniciales $\eta_k(0) = \boldsymbol{\phi}_k^\top
\mathbf{M}\mathbf{q}(0)$ y $\dot{\eta}_k(0) = \boldsymbol{\phi}_k^\top
\mathbf{M}\dot{\mathbf{q}}(0)$. La respuesta temporal completa resulta
de sustituir \eqref{eq:duhamel} en \eqref{eq:modal_exp}:

\begin{equation}\label{eq:resp_completa}
    \mathbf{q}(t)
    = \sum_{k=1}^{3NM}
      \left[
          \eta_k(0)\cos(\omega_k t)
          + \frac{\dot{\eta}_k(0)}{\omega_k}\sin(\omega_k t)
          + \frac{1}{\omega_k}
            \int_0^t f_k(\tau)
            \sin\!\bigl(\omega_k(t-\tau)\bigr)\,\mathrm{d}\tau
      \right]\boldsymbol{\phi}_k.
\end{equation}

Para el caso particular de un impulso $f(t)=F_0\,\delta(t)$ con la
telaraña en reposo ($\mathbf{q}(0)=\dot{\mathbf{q}}(0)=\mathbf{0}$),
la respuesta impulsional es

\begin{equation}\label{eq:impulsional}
    \mathbf{q}(t)
    = F_0\sum_{k=1}^{3NM}
      \frac{\boldsymbol{\phi}_k^\top\mathbf{e}_{n^*,m^*,\hat{d}}}{\omega_k}
      \sin(\omega_k t)\,\boldsymbol{\phi}_k.
\end{equation}

\subsection{Función de transferencia y transmisión entre nodos}

Para cuantificar cuánto de la perturbación aplicada en el nodo fuente
$(n^*,m^*)$ llega al nodo observador $(n^\dagger,m^\dagger)$, se trabaja
en el dominio de la frecuencia. Imponiendo $\mathbf{q}(t)=
\hat{\mathbf{q}}\,e^{i\omega t}$ y $\mathbf{f}(t)=\hat{\mathbf{f}}\,
e^{i\omega t}$ en \eqref{eq:sistema}:

\begin{equation}\label{eq:freq_domain}
    \bigl(\mathbf{K} - \omega^2\mathbf{M}\bigr)\hat{\mathbf{q}}
    = \hat{\mathbf{f}}.
\end{equation}

La solución define la \emph{matriz de receptancia} \cite{ewins2000,inman2014}:

\begin{equation}\label{eq:receptancia}
    \mathbf{H}(\omega)
    = \bigl(\mathbf{K} - \omega^2\mathbf{M}\bigr)^{-1}
    = \sum_{k=1}^{3NM}
      \frac{\boldsymbol{\phi}_k\,\boldsymbol{\phi}_k^\top}
           {\omega_k^2 - \omega^2},
\end{equation}

donde la segunda igualdad es la representación espectral en términos de
los modos normales. La \emph{función de transferencia} entre el nodo
fuente en dirección $\hat{\mathbf{d}}$ y el nodo observador en dirección
$\hat{\mathbf{s}}$ es el escalar

\begin{equation}\label{eq:transferencia}
    H_{obs,src}(\omega)
    = \hat{\mathbf{s}}^\top\mathbf{H}(\omega)\,
      \mathbf{e}_{n^*,m^*,\hat{d}}
    = \sum_{k=1}^{3NM}
      \frac{
          \bigl(\hat{\mathbf{s}}^\top\boldsymbol{\phi}_k\bigr)
          \bigl(\boldsymbol{\phi}_k^\top\mathbf{e}_{n^*,m^*,\hat{d}}\bigr)
      }{\omega_k^2 - \omega^2}.
\end{equation}

El módulo $|H_{obs,src}(\omega)|$ cuantifica la amplitud de transmisión
entre los dos nodos a cada frecuencia. Para comparar configuraciones
geométricas distintas se define el \emph{índice de transmisión integrado}
en la banda $[\omega_a,\omega_b]$:

\begin{equation}\label{eq:indice}
    \mathcal{T}_{obs,src}
    = \int_{\omega_a}^{\omega_b}
      \lvert H_{obs,src}(\omega)\rvert^2\,\mathrm{d}\omega.
\end{equation}

Un valor alto de $\mathcal{T}_{obs,src}$ indica que la configuración
favorece la transmisión en esa banda; uno bajo indica atenuación
\cite{ewins2000,lott2022,mortimer2016}. La comparación de este índice
entre distintos valores de $N$, $M$, $\{r_m\}$, $\{K_m,k_m\}$ y
$\{T^R_m,T^S_m\}$ constituye la herramienta central para responder la
pregunta del proyecto.

\section{Cronograma}
\definecolor{az1bg}{RGB}{230,241,251}   % ciclo 1 fila par
\definecolor{az1alt}{RGB}{238,245,253}  % ciclo 1 fila impar
\definecolor{az1txt}{RGB}{12,68,124}    % ciclo 1 texto ciclo
\definecolor{az1bar}{RGB}{24,95,165}    % ciclo 1 barra lateral
\definecolor{az2bg}{RGB}{234,243,251}   % ciclo 2 fila par
\definecolor{az2alt}{RGB}{242,248,253}  % ciclo 2 fila impar
\definecolor{az2txt}{RGB}{24,95,165}    % ciclo 2 texto ciclo
\definecolor{az2bar}{RGB}{55,138,221}   % ciclo 2 barra lateral
\definecolor{az3bg}{RGB}{220,234,248}   % ciclo 3 fila par
\definecolor{az3alt}{RGB}{232,241,250}  % ciclo 3 fila impar
\definecolor{az3txt}{RGB}{4,44,83}      % ciclo 3 texto ciclo
\definecolor{az3bar}{RGB}{12,68,124}    % ciclo 3 barra lateral
\definecolor{cobg}{RGB}{250,236,231}    % cierre fila par
\definecolor{coalt}{RGB}{250,238,218}   % cierre fila impar
\definecolor{cotxt}{RGB}{153,60,29}     % cierre texto
\definecolor{cobar}{RGB}{216,90,48}     % cierre barra lateral
\definecolor{entbg}{RGB}{252,235,235}   % entregable fondo
\definecolor{entbdr}{RGB}{226,75,74}    % entregable borde
\definecolor{enttxt}{RGB}{121,31,31}    % entregable texto
\definecolor{hdrbg}{RGB}{4,44,83}       % encabezado

\small
\setlength{\tabcolsep}{0pt}
\renewcommand{\arraystretch}{1.5}


\begin{longtable}{
    @{}p{3pt}
    >{\centering\arraybackslash}p{1.1cm}
    p{3.4cm}
    p{7.5cm}
    >{\centering\arraybackslash}p{2.1cm}@{}
}

% ---- encabezado primera página ----
\rowcolor{hdrbg}
& \textcolor{white}{\footnotesize\textbf{Sem.}}
& \textcolor{white}{\footnotesize\textbf{Ciclo}}
& \textcolor{white}{\footnotesize\textbf{Actividad}}
& \textcolor{white}{\footnotesize\textbf{Entregable}} \\
\endfirsthead
\rowcolor{hdrbg}
& \textcolor{white}{\footnotesize\textbf{Sem.}}
& \textcolor{white}{\footnotesize\textbf{Ciclo}}
& \textcolor{white}{\footnotesize\textbf{Actividad}}
& \textcolor{white}{\footnotesize\textbf{Entregable}} \\
\endhead
\endfoot
\endlastfoot


\cellcolor{az1bar} & \cellcolor{az1bg}\textcolor{az1bar}{\textbf{1}}
& \cellcolor{az1bg}\textcolor{az1txt}{\footnotesize Ciclo 1 — Definición}
& \cellcolor{az1bg}{\footnotesize Revisión bibliográfica sobre telarañas orbiculares, dinámica vibratoria y sistemas mecánicos discretos.}
& \cellcolor{az1bg} \\

\cellcolor{az1bar} & \cellcolor{az1alt}\textcolor{az1bar}{\textbf{2}}
& \cellcolor{az1alt}\textcolor{az1txt}{\footnotesize Ciclo 1 — Definición}
& \cellcolor{az1alt}{\footnotesize Configuración geométrica inicial: distribución de nodos, elementos estructurales y variables de interés.}
& \cellcolor{az1alt} \\

\cellcolor{az1bar} & \cellcolor{az1bg}\textcolor{az1bar}{\textbf{3}}
& \cellcolor{az1bg}\textcolor{az1txt}{\footnotesize Ciclo 1 — Modelado}
& \cellcolor{az1bg}{\footnotesize Representación 3D. Grados de libertad y coordenadas generalizadas $\mathbf{q}$.}
& \cellcolor{az1bg} \\

\cellcolor{az1bar} & \cellcolor{az1alt}\textcolor{az1bar}{\textbf{4}}
& \cellcolor{az1alt}\textcolor{az1txt}{\footnotesize Ciclo 1 — Modelado}
& \cellcolor{az1alt}{\footnotesize Formulación de $T$ y $V$ con pretensión. Derivación de la matriz elemental $\mathbf{k}^e_{ij}$.}
& \cellcolor{az1alt} \\

\cellcolor{az1bar} & \cellcolor{az1bg}\textcolor{az1bar}{\textbf{5}}
& \cellcolor{az1bg}\textcolor{az1txt}{\footnotesize Ciclo 1 — Modelado}
& \cellcolor{az1bg}{\footnotesize Ensamblaje de $\mathbf{M}$ y $\mathbf{K}$; verificación del modelo lagrangiano.}
& \cellcolor{az1bg} \\

\cellcolor{az1bar} & \cellcolor{az1alt}\textcolor{az1bar}{\textbf{6}}
& \cellcolor{az1alt}\textcolor{az1txt}{\footnotesize Ciclo 1 — Cierre}
& \cellcolor{az1alt}{\footnotesize Consolidación del modelo, estrategia de solución y documentación.}
& \cellcolor{entbg}\textcolor{enttxt}{\footnotesize\textbf{Anteproyecto}} \\[3pt]


\cellcolor{az2bar} & \cellcolor{az2bg}\textcolor{az2bar}{\textbf{7}}
& \cellcolor{az2bg}\textcolor{az2txt}{\footnotesize Ciclo 2 — Caracterización}
& \cellcolor{az2bg}{\footnotesize Implementación numérica: ensamblaje de $\mathbf{M}$ y $\mathbf{K}$ y verificación de bloques.}
& \cellcolor{az2bg} \\

\cellcolor{az2bar} & \cellcolor{az2alt}\textcolor{az2bar}{\textbf{8}}
& \cellcolor{az2alt}\textcolor{az2txt}{\footnotesize Ciclo 2 — Caracterización}
& \cellcolor{az2alt}{\footnotesize Resolución del problema de valores propios; obtención de $\{\omega_k,\boldsymbol{\phi}_k\}$.}
& \cellcolor{az2alt} \\

\cellcolor{az2bar} & \cellcolor{az2bg}\textcolor{az2bar}{\textbf{9}}
& \cellcolor{az2bg}\textcolor{az2txt}{\footnotesize Ciclo 2 — Caracterización}
& \cellcolor{az2bg}{\footnotesize Clasificación de modos en el plano y transversales; análisis espectral por familia.}
& \cellcolor{az2bg} \\

\cellcolor{az2bar} & \cellcolor{az2alt}\textcolor{az2bar}{\textbf{10}}
& \cellcolor{az2alt}\textcolor{az2txt}{\footnotesize Ciclo 2 — Caracterización}
& \cellcolor{az2alt}{\footnotesize Integral de Duhamel; respuesta temporal ante perturbación impulsional.}
& \cellcolor{az2alt} \\

\cellcolor{az2bar} & \cellcolor{az2bg}\textcolor{az2bar}{\textbf{11}}
& \cellcolor{az2bg}\textcolor{az2txt}{\footnotesize Ciclo 2 — Caracterización}
& \cellcolor{az2bg}{\footnotesize Cálculo de $H_{obs,src}(\omega)$ e índice $\mathcal{T}_{obs,src}$ para la configuración base.}
& \cellcolor{az2bg} \\

\cellcolor{az2bar} & \cellcolor{az2alt}\textcolor{az2bar}{\textbf{12}}
& \cellcolor{az2alt}\textcolor{az2txt}{\footnotesize Ciclo 2 — Cierre}
& \cellcolor{az2alt}{\footnotesize Consolidación de resultados y evaluación del desarrollo del modelo.}
& \cellcolor{entbg}\textcolor{enttxt}{\footnotesize\textbf{Informe de avance}} \\[3pt]


\cellcolor{az3bar} & \cellcolor{az3bg}\textcolor{az3bar}{\textbf{13}}
& \cellcolor{az3bg}\textcolor{az3txt}{\footnotesize Ciclo 3 — Simulación}
& \cellcolor{az3bg}{\footnotesize Definición de configuraciones geométricas para estudio comparativo (variación de $N$, $M$, $r_m$).}
& \cellcolor{az3bg} \\

\cellcolor{az3bar} & \cellcolor{az3alt}\textcolor{az3bar}{\textbf{14}}
& \cellcolor{az3alt}\textcolor{az3txt}{\footnotesize Ciclo 3 — Simulación}
& \cellcolor{az3alt}{\footnotesize Implementación de cada configuración y definición de casos de perturbación mecánica.}
& \cellcolor{az3alt} \\

\cellcolor{az3bar} & \cellcolor{az3bg}\textcolor{az3bar}{\textbf{15}}
& \cellcolor{az3bg}\textcolor{az3txt}{\footnotesize Ciclo 3 — Simulación}
& \cellcolor{az3bg}{\footnotesize Simulación numérica de la respuesta vibratoria y cálculo de $\mathcal{T}_{obs,src}$ por configuración.}
& \cellcolor{az3bg} \\

\cellcolor{az3bar} & \cellcolor{az3alt}\textcolor{az3bar}{\textbf{16}}
& \cellcolor{az3alt}\textcolor{az3txt}{\footnotesize Ciclo 3 — Simulación}
& \cellcolor{az3alt}{\footnotesize Comparación de $\{\omega_k\}$, $\{\boldsymbol{\phi}_k\}$ y $H_{obs,src}(\omega)$ entre configuraciones.}
& \cellcolor{az3alt} \\[3pt]

\cellcolor{cobar} & \cellcolor{cobg}\textcolor{cobar}{\textbf{17}}
& \cellcolor{cobg}\textcolor{cotxt}{\footnotesize Cierre del proyecto}
& \cellcolor{cobg}{\footnotesize Análisis global de resultados, influencia de la geometría y elaboración de conclusiones.}
& \cellcolor{entbg}\textcolor{enttxt}{\footnotesize\textbf{Informe final}} \\

\cellcolor{cobar} & \cellcolor{coalt}\textcolor{cobar}{\textbf{18}}
& \cellcolor{coalt}\textcolor{cotxt}{\footnotesize Cierre del proyecto}
& \cellcolor{coalt}{\footnotesize Artículo científico y preparación de la sustentación.}
& \cellcolor{entbg}\textcolor{enttxt}{\footnotesize\textbf{Sustentación}} \\

\end{longtable}


\printbibliography
\end{document}